\section{Introduction}
\label{sec:intro}

\IEEEPARstart{F}{orecasting} the next few days of a multivariate sensor series supports load planning in power networks, scheduling in industrial plants and maintenance of equipment such as power transformers. Deep forecasters for this task have changed quickly over the last five years. Informer~\cite{informer} and its successors Autoformer~\cite{autoformer} and FEDformer~\cite{fedformer} adapted the Transformer~\cite{vaswani} to long sequences. DLinear~\cite{dlinear} then showed that a single linear layer could match or beat these models on the standard long-term benchmarks, and PatchTST~\cite{patchtst} restored the lead of attention-based models by splitting each channel into short overlapping patches, treating each patch as a token and processing channels independently with shared weights. PatchTST is now a common baseline in forecasting studies and in benchmark libraries~\cite{tslib}.

These models are compared on clean benchmarks, while real sensor streams are noisy in several ways: readings go missing when a logger fails, analog channels pick up noise, transient faults produce isolated spikes, and recalibration or a change in operating regime moves a stretch of the series to a different level or scale~\cite{blazquez_anomaly,gama_drift}. In forecasting there are rarely labels that say which inputs are bad~\cite{seer}, so the model has to cope with corrupted inputs without being told where they are.

Two properties of PatchTST make it sensitive to such inputs. First, every patch enters the encoder with the same standing. The attention weights can in principle learn to ignore a damaged patch, but nothing in the architecture or the training signal encourages this. Second, PatchTST normalizes each input window by its own mean and standard deviation (the instance normalization of RevIN~\cite{revin}, without the affine part) and restores those statistics at the output. This helps under slow drift, but it also removes the information that would tell the model that a window sits at an unusual level relative to the training data. Recent analyses argue that window-local statistics cannot capture every kind of shift~\cite{fan,apt,revinrole}.

In our own baseline run, under the segment-level shift corruption described in Section~\ref{sec:exp}, the test MSE of a faithful PatchTST re-implementation on ETTh1 rises from 0.379 on clean inputs to 0.520 on average over four severity levels, and to 0.712 at the highest level. Spikes raise it to 0.458. Missing values and noise, by contrast, cost less than 0.015.

Existing remedies fall into three groups (Section~\ref{sec:lit}). Normalization-based methods such as RevIN~\cite{revin}, Dish-TS~\cite{dishts}, FAN~\cite{fan} and APT~\cite{apt} target distribution shift but do not weigh inputs by quality. Robustness-oriented frameworks such as SEER~\cite{seer}, RobustTSF~\cite{robusttsf} and TameR~\cite{tamer} address noisy or anomalous inputs, but either replace the forecaster with a new architecture, focus on one type of damage, or act only at training time~\cite{tps}. Expert-based and test-time methods such as TFPS~\cite{tfps} and SOLID~\cite{solid} add parameters or per-sample fine-tuning. None of the studies we surveyed offers a small plug-in for PatchTST that scores the reliability of each patch and gives the head a global view of shift within a single forward pass; we make no claim about work outside that set.

This paper studies whether such a plug-in helps, and by how much. We make four contributions.
\begin{enumerate}
\item A reproducible corruption benchmark on ETTh1 with 22 fixed test conditions (clean plus five corruption families at four severity levels and one mixed setting), in which every model sees identical corrupted inputs.
\item RA-PatchTST, which adds to PatchTST a patch reliability gate, two shift features appended to the forecasting head, and corruption-aware training, for 993 extra parameters in total.
\item An ablation over five variants (A0 to A4) and two seeds that separates the contribution of training on corrupted data from that of each architectural component, together with a direct check of whether the gate down-weights corrupted patches.
\item A browser-based live laboratory in which the original and the proposed model run side by side on real ETTh1 test windows that the user can damage interactively.
\end{enumerate}

RA-PatchTST reduces the average corrupted-input MSE by 3.7\% and the shift error by 11.0\%, but most of this gain comes from corruption-aware training, the gate is a weak detector except for spikes, and the model is slightly worse than PatchTST on clean data and under missing values. The PatchTST backbone is a re-implementation of the published design~\cite{patchtst,patchtst_code}; the research contribution lies in the gate, the shift features, their combination with corruption-aware training, and the ablation that measures what each part adds.

The rest of the paper is organized as follows. Section~\ref{sec:lit} reviews related work. Section~\ref{sec:problem} states the problem formally. Section~\ref{sec:method} develops the method, Section~\ref{sec:arch} describes the system architecture and Section~\ref{sec:impl} the implementation. Section~\ref{sec:exp} gives the experimental protocol and Section~\ref{sec:results} the results and their discussion. Sections~\ref{sec:future} and~\ref{sec:conclusion} outline future work and conclude.
